\begin{proposition}[3.3.4]
\label{IV.2.3.3.4-fr}
Sean $X$, $Y$ dos preesquemas localmente noetherianos, $f:X\to Y$ un
morfismo, $\mathcal E$ un $\mathcal O_X$-módulo coherente, $\mathcal F$ un
$\mathcal O_X$-módulo coherente y plano respecto de $f$. Entonces las
condiciones siguientes son equivalentes~:
\begin{enumerate}
\item[\emph{a)}] $\mathcal E\otimes_Y\mathcal F$ es reducido \textup{(3.2.2)}.
\item[\emph{b)}] Para todo punto
$y\in\operatorname{Ass}(\mathcal E)\cap f(\operatorname{Supp}(\mathcal F))$,
$\overline{\{y\}}$ es un ciclo primo asociado no inmerso de $\mathcal E$,
$\operatorname{long}(\mathcal E_y)=1$ y $\mathcal F_y$ es reducido.
\end{enumerate}
\end{proposition}

Supongamos que se verifica \emph{a)}. Ya se sabe (3.3.3) que, para todo
$y\in\operatorname{Ass}(\mathcal E)\cap f(\operatorname{Supp}(\mathcal F))$,
$\overline{\{y\}}$ es un ciclo primo asociado no inmerso de $\mathcal E$ y
$\mathcal F_y$ carece de ciclo primo asociado inmerso. Además (2.5.5), para
todo $x\in\operatorname{Ass}(\mathcal E\otimes_Y\mathcal F)\cap f^{-1}(y)$,
se tiene
$1=\operatorname{long}((\mathcal E\otimes_Y\mathcal F)_x)
=\operatorname{long}(\mathcal E_y).\operatorname{long}((\mathcal F_y)_x)$,
por tanto $\operatorname{long}(\mathcal E_y)
=\operatorname{long}((\mathcal F_y)_x)=1$, lo que demuestra \emph{b)}.

Recíprocamente, supongamos que se verifica \emph{b)}~; ya se sabe que todo
punto $x\in\operatorname{Ass}(\mathcal E\otimes_Y\mathcal F)$ es un punto
maximal de $\operatorname{Supp}(\mathcal E\otimes_Y\mathcal F)$, que
$y=f(x)$ es un punto maximal de $\operatorname{Supp}(\mathcal E)$ y que $x$
es un punto maximal de $\operatorname{Supp}(\mathcal F_y)$ (3.3.1 y 3.3.3)~;
además, de la hipótesis y de (2.5.5) resulta que
$\operatorname{long}((\mathcal E\otimes_Y\mathcal F)_x)=1$, lo que demuestra
\emph{a)}.

\begin{corollary}[3.3.5]
\label{IV.2.3.3.5-fr}
Sean $X$, $Y$ dos preesquemas localmente noetherianos, $f:X\to Y$ un
morfismo plano~; si $Y$ es reducido en los puntos de $f(X)$ y si $f^{-1}(y)$
es un $\boldsymbol{k}(y)$-preesquema reducido para todo $y\in f(X)$, entonces $X$ es
reducido.
\end{corollary}

Como el nilradical $\mathcal N_Y$ es coherente, el conjunto de los puntos en
los que $Y$ es reducido es \emph{abierto} $(\mathbf{0}_{\mathrm I},5.2.2)$, y es posible
limitarse al caso en que $Y$ es reducido. Basta entonces aplicar (3.3.4) a
$\mathcal E=\mathcal O_Y$ y $\mathcal F=\mathcal O_Y$.
